% !TeX program = xelatex
% Overleaf: Compiler = XeLaTeX; Main document = main.tex.
% 日常写作：直接编译本文件。格式依据：学校 2025 年 10 月指南。
% 隐名辅助版：编译 blind.tex；纸质打印版：编译 print.tex。
% 可选参数：fontset=auto/open/school；bibbackend=bibtex/biber。
\documentclass{tongji-env-master}
% 中英文封面信息。专业字段按本人学籍及入学当年培养方案填写。
% 张三、诸葛亮和 0000000 为示例信息。
\renewcommand\ctitle{市政污水处理工艺优化研究}
\renewcommand\etitle{Optimization of Municipal Wastewater Treatment Processes}
\renewcommand\cauthor{张三}
\renewcommand\eauthor{San Zhang}
\renewcommand\studentid{0000000}
\renewcommand\estudentid{\studentid}
\renewcommand\cschool{环境科学与工程学院}
\renewcommand\eschool{College of Environmental Science and Engineering}
\renewcommand\ccategory{工学}
\renewcommand\ecategory{Engineering}
\renewcommand\cdegree{资源与环境}
\renewcommand\cdegreeen{Resources and Environment}
\renewcommand\edegree{Master of Resources and Environment}
% 2025 级资源与环境培养方案列有“环境工程”领域及“市政工程”研究方向。
% 其他年级应使用其适用培养方案；专业领域与研究方向分别填写。
\renewcommand\cfield{环境工程}
\renewcommand\efield{Environmental Engineering}
\renewcommand\cdirection{市政工程}
\renewcommand\edirection{Municipal Engineering}
\renewcommand\csupervisor{诸葛亮\quad 教授}
\renewcommand\esupervisor{Professor Liang Zhuge}
% 经批准、备案的行业导师限一位；空值自动隐藏该行。
\renewcommand\cindustry{}
\renewcommand\eindustry{}
\renewcommand\cjoint{}
\renewcommand\ejoint{}
% 使用论文成文打印日期；中文日期用汉字。
\renewcommand\cdate{二〇二六年十月}
\renewcommand\edate{October 2026}
% 可选资助信息，建议不超过一项。
\renewcommand\cfunding{}
\renewcommand\efunding{}

\begin{document}
\makecovers
\frontmatter
% 摘要应在主要研究完成后撰写；具体字数按学院及导师要求。
% 下列说明采用正式模板语体。完成摘要后直接以正文替换，无需保留说明。
\begin{cnabstract}
摘要应概括论文的研究目的、研究对象与方法、主要结果及结论，突出本人完成的工作及其创新性。研究结果应以实际数据和分析为依据，结论应明确适用条件与工程意义。

摘要应具有独立性和完整性，语言准确、简练，不列图表。中英文摘要的研究内容、数据、单位及结论应相互对应。关键词一般为三至五个，中文关键词之间用中文逗号分隔。
\keywords{关键词一，关键词二，关键词三}
\end{cnabstract}

% 英文摘要须与最终中文摘要一致；缩写首次出现时给出全称。
\begin{enabstract}
The abstract should concisely present the research objectives, materials and methods, principal results, and conclusions. It should identify the original contribution and its engineering significance, with findings supported by the research evidence.

The abstract should be self-contained and should not include figures or tables. The Chinese and English abstracts should be consistent in their scope, numerical results, units, and conclusions.
\keywordsen{keyword one, keyword two, keyword three}
\end{enabstract}

\tableofcontents
% 可选部分：仅保留正文实际使用的符号；数量较少时可注释 main.tex 的输入行。
\tjunnumbered{符号说明}
\begingroup\tjsmall
\begin{longtable}{@{}p{.17\textwidth}p{.50\textwidth}p{.24\textwidth}@{}}
\toprule
符号 & 含义 & 单位\\\midrule
\endhead
$C_{\mathrm{in}}$ & 进水目标物质量浓度 & \unit{\milli\gram\per\litre}\\
$C_{\mathrm{out}}$ & 出水目标物质量浓度 & \unit{\milli\gram\per\litre}\\
$Q$ & 体积流量 & \unit{\metre\cubed\per\day}\\
$V$ & 体积 & \unit{\metre\cubed}\\
$t$ & 时间 & \unit{\day}\\
$\eta$ & 去除率 & \unit{\percent}\\
\bottomrule
\end{longtable}
\par\clearpage\endgroup
 % 符号较少时可注释。
\mainmatter
% 章节结构可按课题调整，不是学院规定的固定章数。
% 写作建议：先明确具体工程问题，再组织研究现状、问题缺口和研究目标。
\chapter{绪论}
\section{研究背景与工程问题}
本节阐述研究背景、工程需求及研究意义，明确拟解决的问题、研究对象和适用范围。涉及工程规模、水质或运行数据时，应注明来源及时间范围。

污水处理研究涉及水质特征、处理工艺与资源回收等环节\supercite{metcalf2014}。排水系统设计中，工程参数与适用条件应结合相应标准确定\supercite{gb50014}。
% 以上两句展示图书与标准的引用写法，按实际课题替换。
\section{国内外研究现状}
\subsection{相关技术与方法}
本节围绕研究问题梳理代表性文献，比较不同技术与方法的适用条件、主要结论及局限。例如，厌氧氨氧化研究报道了厌氧条件下的铵转化现象\supercite{mulder1995}。
% 该句用于展示期刊论文的引用写法，不限定模板的研究主题。

文献按正文首次出现的顺序编号，引用采用方括号上标，连续引用格式如\supercite{tongji2025guide,tongji2026degree}。
% 上句的学校文件用于展示中文机构作者及网页文献，按正式内容替换。
% 正常引用：研究结论的相应表述\supercite{文献键}。
% 多篇引用：\supercite{key1,key2,key3}；指定页码：\supercite[25--28]{key1}。
% 不手工输入 [1]，不把文献表作为普通文本粘贴。
\subsection{现有研究的不足}
本节概括已有研究尚未解决的问题，说明其与本研究工程对象的联系，明确研究切入点。
\section{研究目标与主要内容}
本节列明研究目标、拟回答的问题及主要研究内容，并说明各项内容之间的逻辑关系。
\section{技术路线与论文结构}
本节说明研究方法、试验验证与工程评价的组织方式，并概述论文各章内容。

% 根据课题选择试验、模型、工程调查等方法，不必保留不适用的小节。
\chapter{研究对象与方法}
\section{研究对象与边界条件}
本节说明工程对象、采样条件、装置规模与运行工况，明确研究边界及试验条件与实际工程条件的关系。
\section{实验设计与分析方法}
\subsection{实验方案与对照设置}
本节说明试验分组、对照条件、平行样数量、运行周期及数据处理方法，并区分独立重复与重复测定。
\subsection{质量保证与质量控制}
本节说明检测依据、仪器设备、校准方法及质量控制措施，必要时列明检出限、定量限和数据有效性判据。
\section{指标计算与单位表达}
物理量采用斜体，单位及描述性下标采用正体。公式按章连续编号，编号置于右侧。以浓度为依据计算去除率时，可表示为
\begin{equation}
  \eta=\frac{C_{\mathrm{in}}-C_{\mathrm{out}}}{C_{\mathrm{in}}}\times100\%
  \label{eq:removal}
\end{equation}
式中，$C_{\mathrm{in}}$与$C_{\mathrm{out}}$分别为进、出水目标物质量浓度，单位为\unit{\milli\gram\per\litre}。使用该指标时，应说明采样条件及浓度比较的适用前提。
% 流量变化明显时需考虑质量通量，不能把浓度去除率直接当作质量去除率。
% 常用写法：\qty{20}{\degreeCelsius}、\unit{\kilo\watt\hour\per\metre\cubed}。
% 化学式：\ce{NH4+}、\ce{NO3-}；pH 正体：\(\mathrm{pH}\)。
\subsection{统计分析与不确定性}
本节说明统计方法、样本量、误差表示方式、模型假设及不确定性来源。

% 表内破折号只表示模板尚未填入数据，不表示零、未检出或缺失观测。
% 定稿时应在表注中明确定义实际采用的缺失值/未检出标记。
\chapter{研究结果与讨论}
\section{运行表现与处理效果}
本节按研究目标呈现主要结果，说明数据范围、运行阶段及评价依据。图表应在正文首次提及后编排，内容与正文相互补充。主要指标可按表\ref{tab:results}组织。
\begin{table}[htbp]
\caption{主要指标记录表}\label{tab:results}
\centering\tjsmall
\begin{tabularx}{\textwidth}{@{}lXXX@{}}
\toprule
指标 & 进水 & 出水 & 评价说明\\\midrule
COD / (\unit{\milli\gram\per\litre}) & — & — & —\\
TN / (\unit{\milli\gram\per\litre}) & — & — & —\\
\bottomrule
\end{tabularx}
\end{table}
\section{影响因素与机制讨论}
本节结合对照试验、物质平衡、表征或模型结果解释主要现象，讨论影响因素及可能机制。研究装置或处理流程可按图\ref{fig:route}组织。
\begin{figure}[htbp]
\centering
% 正式图片写法：\includegraphics[width=.85\textwidth]{figures/文件名.pdf}
\fbox{\parbox[c][28mm][c]{.85\textwidth}{\centering\normalsize 研究装置或工艺流程图}}
\caption{工艺流程图}\label{fig:route}
\end{figure}
\section{本章小结}
本节归纳本章主要发现，说明结论的适用条件及其与后续工程评价的联系。

% 本章为专业硕士的可选组织方式，可按课题改为工艺优化、模型验证、工程设计等。
\chapter{工程应用与综合评价}
\section{工程适用条件与放大约束}
本节分析研究成果的工程适用条件、规模放大要求及运行维护条件，明确试验验证范围与实际应用范围的关系。
\section{技术经济与环境影响评价}
本节说明技术经济或环境影响评价的边界、指标与数据来源，并在一致的比较条件下分析不同方案。
% 如涉及价格，记录价格年份；若课题未开展该项评价，可调整本节。
\section{运行建议与风险控制}
本节提出由试验结果或工程分析支持的运行建议，说明异常工况下的适用限制与控制措施。
\section{本章小结}
本节概括工程应用价值及主要约束，说明研究结果对工程决策的支持程度。

% 结论逐项对应研究目标，不首次引入未经正文论证的新数据。
\chapter{结论与展望}
\section{主要结论}
本节以研究结果为依据，概括论文的主要结论、创新性工作及工程意义，明确结论的适用对象和条件。
\section{研究局限与展望}
本节说明研究的主要局限、尚待解决的问题及后续研究方向。

% 文献按首次引用顺序自动生成；正文使用 \supercite{文献键}。
% 示例文献随正式研究内容替换；不使用 \nocite{*} 列出未引用文献。
% 若开始写作时暂不保留任何引用，可先注释下一行，添加真实引用后恢复。
\printreferences
\startappendices
% 附录为可选部分；不需要时可同时注释 main.tex 的 startappendices 和 input 两行。
% 编号由类文件自动生成：附录 A、B……；公式、图、表 A1、A2……，不手写编号。
\chapter{补充材料}
附录用于列示补充试验记录、计算过程、参数表或其他与论文相关的材料。各附录应分别设置标题，并在正文中给出必要的引用。
\begin{equation}
 V=Qt
 \label{eq:appendix-volume}
\end{equation}
式中，$V$、$Q$与$t$分别为体积、体积流量和时间，其单位应保持一致。
% 下表和图框用于提供附录格式；实际无对应材料时删除。
\begin{table}[htbp]
\caption{补充记录表}\label{tab:appendix}
\centering\tjsmall
\begin{tabularx}{\textwidth}{@{}lXX@{}}\toprule
记录编号&记录内容&来源\\\midrule
—&—&—\\\bottomrule
\end{tabularx}
\end{table}
\begin{figure}[htbp]
\centering\fbox{\parbox[c][15mm][c]{.7\textwidth}{\centering 补充图}}
\caption{补充图}\label{fig:appendix}
\end{figure}
 % 不需要附录时，以上两行一并注释。
\makebackmatter
\end{document}
